\chapter{变换与对称：从图形动作走向群}

\section{把正方形放回框架里}

过年的时候，很多地方还保留着剪窗花的习惯。把一张红纸对折，再对折，拿剪刀随手剪出几个缺口，展开之后，纸上会出现一幅惊人的图案：左边有的花样，右边一定有；上面有的花样，转个方向又出现一次。剪的人根本没有设计过这些重复——它们是对折这个动作“免费赠送”的。

如果你看过万花筒，感受会更直接。几片碎玻璃渣，在三面镜子围成的圆筒里，竟能映出工整得像精密仪器的图案。转动圆筒，图案千变万化，但永远保持着一种说不出的整齐。为什么偏偏是六瓣的花、十二瓣的花，而不是七瓣、十三瓣？因为镜子的夹角是固定的，碎玻璃的影像被镜子反复反射，每反射一次就多复制一份，复制的次数由夹角说了算。玻璃的碎是随机的，图案的整齐却是被镜子的角度强制出来的。

这些“整齐”背后藏着同一个数学问题，而这一章，我们要从一个更朴素的游戏开始研究它。

请准备一张正方形的硬卡片，在它的四个角上分别写上数字 $1$、$2$、$3$、$4$，正反两面都要写，并且保证从任何一面都能看清哪个角是哪个。再在桌面上画一个和卡片一样大的正方形凹槽。现在问：把卡片正好放回凹槽，一共有多少种不同的放法？

注意规则：卡片可以翻面，所以别只想到旋转。

先凭直觉猜一个数。猜 $4$ 的人很多，因为他们只想到了旋转：转 $0^\circ$、$90^\circ$、$180^\circ$、$270^\circ$，正好四种。但卡片还能翻过来，所以答案至少是 $8$。可事情到这里并没有完——真正有意思的问题来了：

如果我把“转 $90^\circ$”和“沿竖直中线翻折”这两个动作接连做，做出来的结果算不算一种“新的放法”？会不会存在第九种放法？再进一步：先旋转再翻折，和先翻折再旋转，是同一件事吗？

这个问题看起来小得像课间游戏，但它难倒过人。十九世纪的数学家们在研究方程时，被迫认真思考“动作接连做”这回事，最后从中提炼出一个全新的数学概念——群。这个概念后来变成了整个现代数学和现代物理的通用语言：晶体的分类、基本粒子的命名、魔方的解法，全都靠它。而这一切，确实可以从桌上这张正方形卡片开始。

\section{先猜，再错}

在发明任何术语之前，我们先看看直觉会把我们带到哪些沟里去。

第一个常见的猜测是：“旋转和翻折接连做，顺序无所谓。”毕竟加法乘法都满足交换律，$3+5=5+3$，凭什么动作不行？这个猜测错得很干脆。拿你的卡片试一下：让角 $1$ 本来在凹槽的左上角。先逆时针转 $90^\circ$，角 $1$ 跑到左下角；再沿竖直中线翻折，角 $1$ 到了右下角。换个顺序：先沿竖直中线翻折，角 $1$ 先到右上角；再逆时针转 $90^\circ$，角 $1$ 落到左上角。两种顺序，角 $1$ 落在完全不同的位置。动作的复合，顺序是要紧的。

\begin{fanli}
警报：“两个动作接连做，先做哪个都一样。”

反例就是上面的卡片实验：先转 $90^\circ$ 再翻折，与先翻折再转 $90^\circ$，角 $1$ 的落点不同。这和数的乘法有本质区别——数的世界里 $3\times5=5\times3$ 天经地义，动作的世界里顺序一变，结果就可能天差地别。生活里你早就知道这一点：先穿鞋再穿袜子，和先穿袜子再穿鞋，绝不是同一回事。
\end{fanli}

第二个常见的猜测是：“对称轴就是横一条、竖一条，共两条。”找正方形的对称轴时，很多同学画出水平中线和竖直中线就停笔了，忘了两条对角线也能让正方形对折后严丝合缝。正方形有四条对称轴，不是两条。少算轴，就会在数“放法”时漏掉一半。

第三个猜测更隐蔽：“既然每个动作都能再做一次、再做一次，那放法应该有无穷多种。”这个猜测错在没注意到“回到原位”这件事。转四次 $90^\circ$，卡片和没转一模一样；翻两次，卡片也恢复原状。动作可以无限地做下去，但做出来的\emph{结果}会不断重复。我们要数的不是“动作序列”——那确实有无穷多条——而是卡片最终的\emph{状态}。把这一点想通，问题才算真正收紧了。

\begin{figure}[htbp]
\begin{center}
\begin{tikzpicture}[scale=1.1]
  \draw[thick] (-2,-2) rectangle (2,2);
  \node at (-2.4, 2.4) {$1$};
  \node at ( 2.4, 2.4) {$2$};
  \node at ( 2.4,-2.4) {$3$};
  \node at (-2.4,-2.4) {$4$};
  \draw[gray, dashed] (-3,0) -- (3,0);
  \draw[gray, dashed] (0,-3) -- (0,3);
  \draw[gray, dashed] (-2.6,-2.6) -- (2.6,2.6);
  \draw[gray, dashed] (-2.6,2.6) -- (2.6,-2.6);
  \draw[-{Stealth}, thick] (50:3.4) arc (50:140:3.4);
  \node at (0,4.1) {$r$};
  \node at (0,-4.3) {四条对称轴（虚线）与逆时针旋转 $90^\circ$ 的动作 $r$};
\end{tikzpicture}
\end{center}
\end{figure}

三次失败给了我们三条线索：动作复合有顺序问题；翻折不止一种；“回到原位”让总数变得有限。要把这些线索织成一张网，我们需要把“动作”本身变成可以计算的数学对象。

\section{把“动作”变成数学对象}

我们平时说“把图形转一下”，说的是一个过程。数学家做的关键一步，是把这个过程看成一个\emph{规则}：平面上的每一个点，都被这个规则搬到了一个新的位置。

\begin{dingyi}[变换]
平面的一个\emph{变换}，是把平面上每个点 $P$ 都搬到某个确定的点 $P'$ 的规则，并且不同的点搬到不同的点，每个点也都被某个点搬到（不丢点、不撞点）。点 $P'$ 叫做点 $P$ 在这个变换下的\emph{像}。
\end{dingyi}

这个定义解决了什么麻烦？它让我们不再需要“盯着图形看”，而可以盯着一个一个的点看。一个变换是否把正方形搬回自身，只要检查正方形的每个点（其实只要检查几个关键点）落到哪里就够了。定义里“不丢点、不撞点”两条限制不是吹毛求疵：如果允许两个点撞到同一个点上，变换就没法撤销——你分不清新来的点原先是谁；撤销不了，后面“逆变换”的故事就讲不下去了。

初中几何里见过的几类变换，都在这里集合了：

\begin{itemize}
  \item \emph{平移}：所有点朝同一个方向移动同样的距离，比如把整张卡片向右推三厘米；
  \item \emph{旋转}：所有点绕某个固定点（旋转中心）转过同一个角度；
  \item \emph{翻折}（也叫轴对称或反射）：所有点翻到某条直线（对称轴）的另一侧，像照镜子；
  \item \emph{缩放}：以某点为中心，所有距离按同一个倍数放大或缩小。
\end{itemize}

前三种有一个共同点：任意两点的距离保持不变。这类变换叫\emph{全等变换}（或等距变换），它们把图形变成和自己全等的图形。缩放就不一样了：距离会变，形状不变而大小变，它把图形变成\emph{相似}图形。这一章的主角是前三种，但缩放在下一章会以新的面目回来。

\begin{dingyi}[复合]
设 $a$ 和 $b$ 是两个变换。先做 $b$、再做 $a$ 得到的整体效果，仍是一个变换，记作 $ab$，叫做 $a$ 与 $b$ 的\emph{复合}。
\end{dingyi}

请特别记住这个约定：$ab$ 是\emph{先做 $b$、后做 $a$}。写成乘法的样子的确方便，但读的时候要“从右往左”读。这个约定来自函数记号——第 2 章里 $f(g(x))$ 也是先算 $g$ 再算 $f$——暂时不习惯没关系，用卡片做两次实验就记住了。

复合的好处立竿见影。比如，两次平移接连做，效果仍是一次平移：先向右推三厘米、再向上推两厘米，合起来就是沿斜方向推一段，距离和方向都能算出来。两次旋转呢？绕\emph{同一个}中心转 $30^\circ$ 再转 $45^\circ$，等于绕这个中心转 $75^\circ$——角度直接相加。可如果两次旋转的中心不同，事情就有趣了：结果仍是一次旋转（或一次平移），但新的旋转中心藏在哪儿，得画辅助线找半天。这正是“动作接连做”值得研究的原因：每一步都简单，合起来却有惊喜。

还有一个最“懒”的变换：什么都不做，每个点原地不动。它叫\emph{恒等变换}，记作 $e$。别小看它，它在整个理论里的地位，相当于乘法里的 $1$。

\begin{dingyi}[对称变换]
如果一个变换把图形 $F$ 的每个点都搬到 $F$ 内部的点上，使整个图形变换后与原来完全重合，就称这个变换是图形 $F$ 的一个\emph{对称变换}。
\end{dingyi}

注意“完全重合”四个字。把正方形转 $45^\circ$ 不是它的对称变换——四个角戳出去了。把正方形向右平移三厘米也不是——图形挪窝了。对称变换是一类特殊的全等变换：它不但保持距离，还保持图形这个整体。

这里有个耐人寻味的差别：正方形只有寥寥几种对称变换，而圆有无穷多种——绕圆心转任何角度，圆都和自身重合，沿任何一条直径翻折也一样。直觉上我们会说“圆比正方形更对称”，现在这句话有了精确的度量：一个图形的对称变换越多，它就越对称。正方形的“对称程度”是一个有限的清单，我们马上就能把它全部清点出来。

现在，正方形卡片的问题有了精确的语言：我们问的其实是“正方形有多少个对称变换”，而“动作接连做”就是对称变换的复合。

\begin{toolbox}
本章的新工具，请收好：
\begin{itemize}
  \item \emph{变换}：把“动一下图形”翻译成“搬运每个点的规则”，从此动作可以精确比较；
  \item \emph{复合} $ab$：先做 $b$ 再做 $a$，把“接连做”变成可以列式计算的对象；
  \item \emph{恒等变换} $e$：什么都不做，是计数和比较的基准；
  \item \emph{逆变换}：把变换的效果撤销。转 $90^\circ$ 的逆是转 $270^\circ$，翻折的逆是它自己；
  \item \emph{对称变换}：让图形与自身重合的变换，是我们清点“整齐程度”的清单条目。
\end{itemize}
\end{toolbox}

\begin{shuyu}{逆变换}
\emph{直观解释}：动作的“撤销键”。\emph{正式定义}：若变换 $a$ 与变换 $b$ 满足 $ab=e$ 且 $ba=e$，就称 $b$ 是 $a$ 的逆变换，记作 $a^{-1}$。\emph{例子}：逆时针转 $90^\circ$ 的逆变换是顺时针转 $90^\circ$（也就是逆时针转 $270^\circ$）；沿任何轴翻折的逆变换是这次翻折本身，因为翻两次就回到了原样。
\end{shuyu}

\section{八种，一种不多}

回到卡片问题。我们现在可以给出完整的答案，并且\emph{证明}它。

用 $r$ 表示绕正方形中心逆时针转 $90^\circ$，用 $f$ 表示沿竖直中线翻折。那么 $r^2=rr$ 是转 $180^\circ$，$r^3$ 是转 $270^\circ$，而 $r^4=e$。看起来我们能造出的动作越来越多，但下面的定理说：全是这八个里头的。

\begin{dingli}[正方形的对称变换]
正方形恰好有 $8$ 个对称变换：
\[
e,\quad r,\quad r^2,\quad r^3,\quad f,\quad rf,\quad r^2f,\quad r^3f,
\]
即四个旋转（含恒等变换）和四个翻折。
\end{dingli}

\begin{proof}
先证明“不超过 $8$ 个”。设 $T$ 是正方形的任意一个对称变换。由于对称变换保持距离，正方形的中心（对角线交点）是唯一到四个顶点距离都相等的点，所以 $T$ 必然把中心搬到中心。于是四个顶点只能被搬到四个顶点。看角 $1$：它的像至多有 $4$ 种选择。角 $1$ 的像确定后，看与角 $1$ 相邻的角 $2$：因为 $T$ 保持距离，角 $2$ 的像必须落在“角 $1$ 的像”的相邻顶点上，而至多只有 $2$ 个相邻顶点。最后，一旦角 $1$ 和角 $2$ 的像都确定了，整个正方形的位置就被唯一确定——由“边边边”全等判定，以这两个相邻顶点为一条边的正方形只有一种摆法，其余两个顶点无处可逃。所以 $T$ 至多有 $4\times2=8$ 种可能。

再证明“这 $8$ 种都能做到”。$e,r,r^2,r^3$ 是四个旋转，显然都是对称变换。$f$ 是沿竖直中线的翻折；$rf$ 是“先沿竖直中线翻折、再转 $90^\circ$”，直接检查顶点落点，它等价于沿一条对角线的翻折；同理 $r^2f$、$r^3f$ 分别等价于沿水平中线和另一条对角线的翻折。正方形恰有四条对称轴，所以这四个翻折都是对称变换，且与四个旋转互不相同。合起来恰好 $8$ 个，一种不多，一种不少。
\end{proof}

\begin{tuilun}
正方形的每个对称变换都有逆变换，而且逆变换仍在这八个之中。
\end{tuilun}

\begin{proof}
旋转 $r$ 的逆是 $r^3$，因为 $rr^3=r^4=e$；$r^2$ 的逆是它自己；四个翻折的逆都是它们自己，因为沿同一条轴翻两次必回原位；$e$ 的逆是 $e$。八个成员，人人有逆，且逆不出这个圈子。
\end{proof}

\begin{zhu}
证明里“由边边边全等判定，正方形只有一种摆法”这一步，其实说出了一个更普遍的事实：平面上的全等变换，被它不共线的三个点的像完全决定。正方形有四个顶点，但我们只需要盯住两个相邻的——因为中心已被锁定，第三点的选择自动收紧。这就是为什么“$4\times2$”这个小小的乘法能管得住整个图形。
\end{zhu}

这个证明的方法值得记住：要数清楚“有多少种动作”，不用把动作一个个表演出来，只需要盯着\emph{关键点的像}。两个相邻顶点的去向，锁定了整个正方形。这种“用少量信息锁定整体”的想法，后面还会反复出现。

八个变换清点了，接下来是最能体现本章精神的一步：把它们的复合关系\emph{全部}算出来，列成一张表。表中第 $x$ 行、第 $y$ 列的格子，写的是先做 $y$、再做 $x$ 的结果 $xy$。

\begin{table}[htbp]
\centering\small
\begin{tabular}{c|cccc|cccc}
\toprule
 & $e$ & $r$ & $r^2$ & $r^3$ & $f$ & $rf$ & $r^2f$ & $r^3f$\\
\midrule
$e$     & $e$ & $r$ & $r^2$ & $r^3$ & $f$ & $rf$ & $r^2f$ & $r^3f$\\
$r$     & $r$ & $r^2$ & $r^3$ & $e$ & $rf$ & $r^2f$ & $r^3f$ & $f$\\
$r^2$   & $r^2$ & $r^3$ & $e$ & $r$ & $r^2f$ & $r^3f$ & $f$ & $rf$\\
$r^3$   & $r^3$ & $e$ & $r$ & $r^2$ & $r^3f$ & $f$ & $rf$ & $r^2f$\\
\midrule
$f$     & $f$ & $r^3f$ & $r^2f$ & $rf$ & $e$ & $r^3$ & $r^2$ & $r$\\
$rf$    & $rf$ & $f$ & $r^3f$ & $r^2f$ & $r$ & $e$ & $r^3$ & $r^2$\\
$r^2f$  & $r^2f$ & $rf$ & $f$ & $r^3f$ & $r^2$ & $r$ & $e$ & $r^3$\\
$r^3f$  & $r^3f$ & $r^2f$ & $rf$ & $f$ & $r^3$ & $r^2$ & $r$ & $e$\\
\bottomrule
\end{tabular}
\end{table}

这张表值得你趴在桌上研究十分钟。先看左上角 $4\times4$ 的小方块：旋转复合旋转还是旋转，而且规律就是“角度相加、满 $360^\circ$ 归零”，所以它其实是加法表的翻版。再看右下角 $4\times4$：翻折复合翻折，得到的却是旋转——照镜子照两次（沿不同的轴），效果等于一个旋转，而且转角恰好是两轴夹角的两倍。这一点你可以拿两面小镜子摆个角度亲自看到。

还有两个重要观察。第一，表里每个格子都还是那八个成员之一：八个动作在复合下是\emph{封闭}的，任做多少步都逃不出去，所以根本不存在“第九种放法”。第二，表关于主对角线\emph{不对称}：比如 $r$ 行 $f$ 列是 $rf$，而 $f$ 行 $r$ 列是 $r^3f$。也就是说
\[
fr \;=\; r^3 f \;\neq\; rf.
\]
这正是第一节卡片实验的代数形式：复合不满足交换律。

\begin{shiyan}
用实物验证表里的一格。拿你的正方形卡片，四角标好 $1,2,3,4$。先沿竖直中线翻折（这就是 $f$），再逆时针转 $90^\circ$（这就是 $r$），记下四个角的位置——你刚才做的是 $rf$。把卡片复位，反过来做：先转 $90^\circ$，再翻折，这是 $fr$。对照复合表，看看 $fr$ 那一格里写的 $r^3f$ 和 $rf$ 是不是不同的角位置。再挑三格验证，直到你对这张表心服口服。
\end{shiyan}

现在退后一步，看这八个动作的\emph{整体}。它们满足四条性质：复合封闭（格子里没有外人）；有恒等元 $e$（和任何动作复合都不改变对方）；每个动作都有逆（每行每列都恰好出现一次 $e$，转 $90^\circ$ 配转 $270^\circ$，翻折配自己）；复合满足结合律（$(ab)c=a(bc)$，因为两边都是“先做 $c$、再做 $b$、再做 $a$”，括号怎么括只是说话的节奏）。

结合律值得多说一句。它听起来像废话，却是整个体系能运转的地基：正因为有它，“先做 $c$、再做 $b$、再做 $a$”这串长动作才能不加分号、不加括号地写成 $abc$，谁都不必争论先算哪一段。对比一下：减法就没有这条性质，$(8-5)-2=1$ 而 $8-(5-2)=5$，所以连减必须老老实实加括号。动作的复合天生自带结合律，这是我们的运气——也是“群”这套语言能如此简洁的原因之一。

满足这四条性质的系统，有一个专门的名字。

\begin{shuyu}{群}
\emph{直观解释}：一袋子可以互相复合、可以撤销、复合结果跑不出袋子的“动作”。\emph{正式定义}：一个集合 $G$ 配上一种运算，满足：运算封闭；有恒等元；每个元素都有逆元；运算满足结合律。这样的 $(G,\text{运算})$ 称为一个\emph{群}。\emph{例子}：正方形的八个对称变换配上复合，构成一个群，叫做正方形的对称群。它有 $8$ 个元素，不满足交换律。
\end{shuyu}

这就是本章标题里那个“群”。它不是天上掉下来的术语，而是被逼出来的：当我们认真追问“动作接连做”的规律时，这四条性质自己浮出了水面。数学家们后来发现，值得研究的对称系统——图形的、方程的、晶体的、粒子的——全都恰好满足这四条。于是群成了研究一切对称的统一语言。

\section{对称到底有什么用}

学会一个新概念的标志，是它能帮你做成以前做不到的事。群和对称至少在三件大事上立下过功劳。

\emph{第一件：给图案分类。}剪窗花时，对折一次剪出的图案有一条对称轴；对折两次再剪，图案通常有两条对称轴，甚至更多。墙纸、地砖、布料上的图案沿着整个平面无限重复，数学家问：平面上的重复图案，按对称变换的组成来分，一共有多少种？答案令人吃惊：不多不少，恰好 $17$ 种。也就是说，从古埃及的壁画到现代的瓷砖，人类能贴出来的所有平面重复图案，无论花纹多花哨，其对称结构必然是这 $17$ 种之一。把问题升到三维——晶体的原子排列——答案变成 $230$ 种，化学家至今仍用这张表来登记晶体。这些结论的证明超出了本书范围，但计数用的工具你已经见过：就是“复合封闭、有逆、有恒等”那一套。

\begin{lianjie}
回头看看：第 3 章讲的全等三角形，本质就是“存在一个全等变换把一个三角形搬到另一个上”；第 5 章讲多项式的韦达定理，马上要在下面重新登场，它是“方程根的对称”的化身。往前看看：第 7 章会告诉你，旋转、翻折、缩放这些变换可以写成数表（矩阵）来做计算。对称不是孤立的景点，它是连通代数、几何和物理的立交桥。
\end{lianjie}

\emph{第二件：解方程。}这件功劳最漂亮，也最出人意料。以二次方程 $x^2-5x+6=0$ 为例，它的两个根是 $2$ 和 $3$。现在做一个“动作”：把两个根互换位置。哪些式子在这个动作下不变？$x_1+x_2=5$ 不变，$x_1x_2=6$ 不变——交换两个数，和与积都不变。而 $x_1-x_2$ 会变号。换句话说，$x_1+x_2$ 和 $x_1x_2$ 是“根互换”这个变换的\emph{对称式}，而韦达定理说的正是：这些对称式可以直接从方程的系数读出来，根本不用先求根。

把这个想法推向三次、四次方程，数学家们发现：方程能不能用根式解出来，取决于它的根的“对称群”的某种内部结构。五次方程的群结构太复杂，于是一般五次方程没有根式解——这个震撼数学界的结论，证明里用到的核心概念，正是我们桌上这八个动作所示范的群。当年创立这套理论的年轻数学家们，是在和当时整个数学界的怀疑作斗争的；他们关于“对称决定可解性”的洞见，后来长成了代数学的主干。

\emph{第三件：驯服魔方。}一个打乱的三阶魔方，所有可能的彩色图案共有 $43{,}252{,}003{,}274{,}489{,}856{,}000$ 种，大约是 $4.3\times10^{19}$，念出来约是“四千三百亿亿”。如果靠瞎转还原，宇宙的年龄都不够用。但魔方的所有转法构成一个群：每一次转动是一个变换，转法可以复合，可以撤销。群论的语言让玩家能说清楚“这一段操作只交换了两个棱块、其余不动”——这类只改动极少块的操作叫做\emph{交换子}，是一切魔方公式的原料。数学家还借助计算机证明：任何打乱状态，至多 $20$ 步（把转 $180^\circ$ 也算一步）就能还原。这个数字不是猜的，是靠群的结构把天文数字的局面先归并成等价类，再逐一核验得到的。

三件事用的是同一把钥匙：不去追踪每一个具体的对象，而去追踪\emph{对象之间的变换}；不满足于“形状相同”，而追问“结构相同”。

什么是“结构相同”？看一个最小的例子。正方形的四个旋转 $e,r,r^2,r^3$ 自己封闭，构成一个小群。另一边，考虑“钟面上只有 $0,1,2,3$ 四个刻度、做模 $4$ 加法”的系统：$1+2=3$，$2+3=1$，$3+3=2$。这两个系统，一个是几何动作，一个是算术，成员完全不同。但把 $r$ 对应到 $1$，$r^2$ 对应到 $2$，$r^3$ 对应到 $3$，$e$ 对应到 $0$，你会发现两边的运算表一模一样——旋转 $90^\circ$ 再转 $180^\circ$ 等于转 $270^\circ$，正如 $1+2\equiv3\pmod 4$。两张表只是标签不同，骨架全同。这时我们说：这两个群\emph{结构相同}（术语叫“同构”）。

这一步跨出去很大。几何动作和整数加法，看上去风马牛不相及，却在结构的层面上是同一个东西。数学的力量正来自这种“换个标签看世界”的本领：只要结构相同，在一侧证明的结论，立刻免费搬运到另一侧。比如模 $4$ 加法里 $1$ 重复加自己能生成所有元素，立刻就知道旋转群里 $r$ 重复复合也能生成全部四个旋转；这种“一个成员生出全家”的群有个专门的名字，叫循环群。反过来，正方形完整的八元素群就找不到这样一个成员——$r$ 生不出翻折，$f$ 翻来覆去只有它和 $e$ 两个。结构的眼光，让“八个动作的大家庭”和“四个旋转的小家庭”的差别一目了然。

回头看全章的旅程：我们从“形状相同”（全等）出发，经过“动作相同”（对称变换），最后抵达“结构相同”（群与同构）。抽象程度一级比一级高，但每一级都不是为了抽象而抽象——站得越高，一次能看清楚的东西就越多。

\begin{pandeng}
想继续往上爬的同学，可以记住这些关键词：\emph{群论}（研究对称结构的分支）、\emph{同构}（结构相同的精确定义）、\emph{伽罗瓦理论}（用群判定方程可否根式解）、\emph{置换群}（魔方的数学模型）、\emph{壁纸群与空间群}（$17$ 种与 $230$ 种的故事）、\emph{李群}（连续的对称，物理学的语言）。这些名词今天只是路标，但它们全都从一张正方形卡片出发。
\end{pandeng}

\section*{思考题}

\paragraph{观察题}
\begin{sikao}
  \item 把本章的正方形换成等边三角形，数一数它有多少个对称变换，并把旋转和翻折分开数。提示：模仿定理的证明——顶点 $1$ 的像有几种选择？相邻顶点的像还剩几种？
  \item 在正方形的复合表中，找出所有满足 $x^2=e$ 的元素，数一数有几个。提示：除了四个翻折，旋转里也藏着一个“做两次回到原位”的成员。
\end{sikao}

\paragraph{证明题}
\begin{sikao}
  \item 证明：若 $a$ 和 $b$ 都是某个图形 $F$ 的对称变换，则复合 $ab$ 也是 $F$ 的对称变换。提示：回到定义，检查“变换后图形与自身重合”这句话经过两步之后是否仍然成立。
  \item 证明脱鞋脱袜定律：$(ab)^{-1}=b^{-1}a^{-1}$。提示：把 $ab$ 与 $b^{-1}a^{-1}$ 复合起来，利用结合律一步步消成 $e$；再想想为什么顺序要反过来——早上先穿袜子后穿鞋，晚上就得先脱鞋后脱袜。
\end{sikao}

\paragraph{挑战题}
\begin{sikao}
  \item 一个魔方只靠六面的转动来玩。问：保持立方体整体重合的\emph{旋转}对称变换有多少种（不许翻折）？提示：盯住一个面——它有几种去处？面定了之后，立方体还能绕哪根轴转？答案是一个让你有点小得意的两位数。
  \item 设计一个只有两个元素的动作系统，使它构成群，并写出它的复合表。提示：从“翻折加恒等变换”或“奇数偶数相加”里挑一个下手，再验证四条群规则。
\end{sikao}

\section*{通往下一章的门}

这一章里，我们把“转一下”“翻一下”变成了能列进表格、能计算的对象。但你可能已经隐隐不满：验证一个复合，要拿卡片摆弄；算一个旋转之后某点落到哪里，要靠画图量取。动作在我们手里还是\emph{手工活}。

有没有办法让变换变成算术？比如，把平面上每个点用两个数 $(x,y)$ 记下位置，那么“绕原点转 $90^\circ$”能不能写成一套对这两个数的固定算法——不管什么点代进去，转完的新坐标立刻算出来？“向右平移三厘米”呢？缩放两倍呢？如果能做到，那么复合变换就是两种算法的接连执行，复合表可以用算式自动生成，卡片可以彻底退休。

可以做到。为此我们需要一个新的主角：一种既有长度又有方向的量——\emph{向量}，以及操纵它的数表——\emph{矩阵}。下一章，二维箭头将带我们走进线性代数的世界，而本章的旋转、翻折和缩放，会在那里变成几行干净利落的数字。
